\chapter{Система SI и Стандартная модель}
\label{ch:si-sm}

\section{Иерархия масштабов}

От ячейки~$\hv$ к первичному дефекту, составным узлам, макроскопическим частицам
и, наконец, полям Стандартной модели --- каждый уровень есть осреднение предыдущего
без новых свободных параметров в~$g$.

\section{Мост к системе SI}

\begin{definition}[Планковская идентификация]
\label{def:si-bridge}
$\ell_P:=\hl$, $a=\hl$, $\kgeom$ из однотактового тела (гл.~\ref{ch:carrier}):
\[
  \htick=\kgeom\, t_P,\quad E_0=\kgeom E_P,\quad c=\kgeom c_0.
\]
\end{definition}

Планковское время $t_P=\ell_P/c$ уже содержит макроскопическую скорость~$c$;
$\kgeom$ выводится из геометрии, после чего $c=\kgeom c_0$ --- проверка согласованности.

\section{Постоянная тонкой структуры}

\begin{theorem}[Константа $\alpha_{\mathrm{fs}}$]
\label{thm:alpha}
\[
  \alpha_{\mathrm{fs}}^{-1}=4\pi^3+\pi^2+\pi\approx 137.036.
\]
Формула --- анзац фазового объёма по степеням $\pi$.
\end{theorem}

\begin{proof}
\emph{Микроуровень.}
Минимальный фазовый шаг на такт
\[
  \Delta\phi_{\min}=\tfrac12.
\]

Полный оборот фазы есть $2\pi$.
Доля одного микрошага в этом обороте
\[
  \frac{\Delta\phi_{\min}}{2\pi}=\frac{1/2}{2\pi}=\frac{1}{4\pi}.
\]
В saturating gate эта доля сидит в числителе при $\rho_E\to 0$:
\[
  \alpha^*-1=\frac{1}{4\pi},\qquad \alpha^*=1+\frac{1}{4\pi}.
\]
Число $\alpha^*$ --- поправка одного такта на~$M$.
Оно не равно $1/137$ и не содержит $|N|$.

\smallskip
\emph{Макроуровень.}
Константа тонкой структуры на~$T$ есть безразмерная сила электромагнитного сопряжения $e^2/(4\pi\varepsilon_0\hbar c)$.
Телесный угол в трёх измерениях есть $4\pi$, и он не равен числу соседей.
Выражение, собранное из этого угла и из плоских фазовых факторов $\pi^2$ и $\pi$,
\[
  \alpha_{\mathrm{fs}}^{-1}=4\pi^3+\pi^2+\pi,
\]
не содержит $|N|$ и не совпадает с $\alpha^*$.
Почленный вывод этих трёх слагаемых в этой главе не закрыт:
формула фиксирует фазовый объём на~$T$ и проверяется числом.
Подстановка $|N|$ вместо $4\pi$ дала бы другое число и на гексагональном срезе, и на FCC.
\end{proof}

\section{Электрослабый сектор}

Пиковая плотностная шкала (\cref{sec:theorem-51}, \cref{sec:bit-budget}):
\[
  N_{\mathrm{hier}}=\lfloor B_{\hv}\rfloor-1=8,\qquad
  v_{\mathrm{peak}}=\alpha_{\mathrm{fs}}^{8} E_P.
\]

\begin{theorem}[Вакуумное среднее значение]
\label{thm:vev}
\[
  v=v_{\mathrm{peak}}\sqrt{2\pi}=\alpha_{\mathrm{fs}}^{8} E_P\sqrt{2\pi}\approx 246\ \mathrm{GeV}.
\]
\end{theorem}

\begin{proof}
По \cref{sec:bit-budget} информационный бюджет ячейки $B_{\hv}=2\pi/\ln 2$,
$N_{\mathrm{ring}}=2^{\lfloor B_{\hv}\rfloor}=512$.
Число иерархических ступеней ослабления на~$\Tlayer$
\[
  N_{\mathrm{hier}}=\lfloor B_{\hv}\rfloor-1=8
\]
--- битовая глубина минус один бит занятости pra-узла (\cref{ch:matter}).

\smallskip
\emph{Пик.}
Каждый электрослабый канал ослабляется на один фактор~$\alpha_{\mathrm{fs}}$
(\cref{thm:alpha}); после $N_{\mathrm{hier}}$ ступеней
\[
  v_{\mathrm{peak}}=\alpha_{\mathrm{fs}}^{\,N_{\mathrm{hier}}} E_P
  =\alpha_{\mathrm{fs}}^{8} E_P.
\]

\smallskip
\emph{Интеграл readout'а.}
Теорема~\ref{thm:t-cr}: осреднение~$\mathcal{B}$ на масштабах $\gg\hl$
даёт гауссово ядро ширины~$\sigma_R$.
Для стандартной нормали отношение пиковой плотности к интегральной амплитуде
равно $\sqrt{2\pi}$:
локальный пик $v_{\mathrm{peak}}$ на~$M$ readout'ится как
\[
  v=v_{\mathrm{peak}}\sqrt{2\pi}\approx 246\ \mathrm{GeV}.
\]
Множитель $\sqrt{2\pi}$ --- не подгонка, а следствие биномиального осреднения
(\cref{def:soft-readout}) в гауссовом пределе.
\end{proof}

\section{Массы частиц}

\begin{proposition}[Лептоны и барионы]
\begin{align*}
  m_H^{(\mathrm{bare})}&=v/2,\\
  m_e&\approx \alpha_{\mathrm{fs}}^2 m_H/N_\phi,\quad N_\phi=13,\\
  m_p&\approx \alpha_{\mathrm{fs}}\, v/2\cdot(1+\kgeom^2/N_{12}),\\
  m_n&=m_p+2m_e\quad\text{(канал $\beta$-распада)}.
\end{align*}
\end{proposition}

Численное сравнение с данными PDG приведено в приложении; здесь зафиксированы формулы вывода.

\section{Четыре взаимодействия}

\begin{proposition}[Один закон --- четыре макроскопических канала]
Электромагнетизм: фазы на рёбрах, $\Phi_\square$, излучение с $n=0$.
Слабое: лево- и право-спиновые проекторы, переворот $SU(2)$; $\sin^2\theta_W=3/13$
из геометрии кластера $12+1$.
Сильное: составные узлы, $\alpha_s(v)=d/(N_{\mathrm{hier}}\pi)$.
Гравитация: деформация метрики на~$T$,
$G_{\mu\nu}=\mathcal{E}(\delta[\mathrm{strain}])=8\pi\ell_P^2 T_{\mu\nu}$.
\end{proposition}

Полный вывод гравитационных волн, полного спектра барионов и матрицы CKM
оставлен для дальнейшей работы и не ослабляет формулировок теорем выше.
